\documentclass{subfiles}
\begin{document}
    Red-Black trees are a special type of \gls{bst},
        from which they differ for the following properties:
        \begin{itemize}
            \item Each node of a RB tree is either red or black;
            \item The root of a RB tree is always black;
            \item Each leaf is black;
            \item Each red node has only black children; and 
            \item For each node, all simple path from the node to a descendant leaf,
                contains the same number of black nodes.
        \end{itemize}
    
    For any RB tree \(T\), the following result holds.
    \begin{lemma*}
        If \(T\) is a Red-Black tree with \(n\) internal nodes, 
            then \(T\) has height at \(2 \log(n + 1)\).        
    \end{lemma*}
    The proof of the above Lemma is based on the notion of black height of a RB tree,
        which essentially represents the number of black nodes in the path from 
        a node (the proof considers the root) to any descendant leaf.

    In \cref{sec:RB-rotations} we describe the rotation operations needed to implement the dictionary  
        operations, later described in \Cref{sec:RB-search,sec:RB-insert,sec:RB-delete}.

    \subsubsection{Rotations}
    \subfile{../subsub/2.1.1 - RB rotations}
\end{document}
